\section{Experiments}
\label{sec:experiments}

\subsection{Fixed-budget protocol}

We study nine D4RL locomotion tasks: HalfCheetah, Hopper, and Walker2d with medium, medium-replay, and expert datasets \citep{fu2020d4rl}. The implementation follows the released JAX TD3+BC sweep code and CORL-style hyperparameters \citep{tarasov2022corl}. The total nominal actor budget is
\begin{equation*}
\begin{gathered}
 T\in\{0.05,0.1,0.2,0.4,0.7,1.5,2.5,\\
 4,7,10,12,14,17,20\}.
\end{gathered}
\end{equation*}
Each main condition uses two seeds, ten deterministic evaluation episodes, and $10^6$ critic updates. The complete matrix now contains six methods: TD3+BC ($K=1$), MPI-Prox for $K\in\{2,3,4\}$, and MPI-Lin for $K\in\{2,3\}$, for $6\times9\times14\times2=1512$ final scores.

We define \emph{high budget} as $T\geq4$ and \emph{tail budget} as $T\geq10$. A high-budget environment--budget cell is descriptively called collapsed when its two-seed mean score is below $20$. We also report raw seeded-run counts and threshold sensitivity. This threshold is an operational locomotion failure marker, not a theorem or a universal offline-RL criterion.

\subsection{Additional hops broaden the observed stability envelope}

\begin{figure*}[t]
\centering
\includegraphics[width=\textwidth]{figures/stability_envelope.pdf}
\caption{Complete fixed-budget sweep. Proximal $K=4$ extends the high-return region and has the lowest high-budget failure rate. Linearized $K=3$ also improves over linearized $K=2$, so the effect is not unique to the proximal loss.}
\label{fig:stability}
\end{figure*}

\begin{table}[t]
\centering
\caption{Aggregate fixed-budget results. Collapse entries are two-seed cells / raw seeded runs for $T\geq4$ and score $<20$.}
\label{tab:aggregate}
\footnotesize
\begin{tabular}{lrrrr}
\toprule
Method & $T\leq2.5$ & $T\geq4$ & $T\geq10$ & Collapse\\
\midrule
TD3+BC ($K=1$) &68.35&25.97&20.46&35/68\\
MPI-Prox ($K=2$)&\textbf{71.62}&42.02&31.49&25/53\\
MPI-Prox ($K=3$)&70.56&52.78&43.55&17/41\\
MPI-Prox ($K=4$)&71.21&\textbf{63.61}&\textbf{56.26}&\textbf{9/29}\\
MPI-Lin ($K=2$)&69.02&36.68&27.44&29/58\\
MPI-Lin ($K=3$)&68.54&43.91&33.12&24/53\\
\bottomrule
\end{tabular}
\end{table}

Figure~\ref{fig:stability} and Table~\ref{tab:aggregate} show a consistent descriptive pattern. The high-budget mean increases from $25.97$ at $K=1$ to $42.02$, $52.78$, and $63.61$ for proximal $K=2,3,4$. Corresponding collapse counts decrease from $35/63$ to $25/63$, $17/63$, and $9/63$. The ordering is unchanged for thresholds $10$, $20$, $30$, and $40$ (Appendix~\ref{app:tables}). The gain is not purchased by a large loss at smaller budgets: the $T\leq2.5$ mean is $71.21$ for $K=4$, compared with $68.35$ for $K=1$.

The new $K=4$ grid sharpens the lower-tail interpretation. Across the 63 paired high-budget environment--budget cells, $K=4$ exceeds $K=3$ by more than one point in 26 cells, lies within one point in 22, and loses by more than one point in 15. The mean difference is $+10.83$, whereas the median is only $+0.32$. At the collapse threshold, $K=4$ rescues eight cells that are collapsed for $K=3$, creates no new collapsed cell, leaves nine cells collapsed under both, and leaves 46 stable under both. These are descriptive paired-cell comparisons; they do not correct the $K$-dependent initialization issue discussed below.

\begin{figure*}[t]
\centering
\includegraphics[width=\textwidth]{figures/k4_analysis.pdf}
\caption{Descriptive $K$ scaling for the proximal family. (a) High-budget mean and collapse rate. (b) Cellwise $K=3$ versus $K=4$ comparison over 63 high-budget cells. The improvement is concentrated in failure rescue rather than a uniform shift.}
\label{fig:k4}
\end{figure*}

The linearized control gives a weaker but directionally similar result. Increasing from $K=2$ to $K=3$ raises the high-budget mean from $36.68$ to $43.91$ and reduces collapsed cells from $29$ to $24$. Therefore, the observed stabilization is not unique to optimizing the proximal loss. Because the two realizations use different actor objectives and both use amortized neural regression, this comparison does not identify an implicit-versus-explicit integration effect.

\subsection{Task dependence and baseline replication}

The high-budget benefit of $K=4$ is heterogeneous (Appendix Table~\ref{tab:env_high}). Relative to $K=1$, it is largest on Hopper-medium, Hopper-expert, Walker2d-medium, and Walker2d-expert; HalfCheetah-medium is essentially unchanged. The per-budget curves are also nonmonotone. For example, $K=4$ Hopper-medium remains near 100 through $T=17$ but becomes seed-unstable at $T=20$, while Walker2d-expert is stable at $T=12$, collapses at $T=14$, and separates sharply across seeds at $T=17$. We therefore avoid a universal claim such as a frontier proportional to $K$.

The one-hop phase diagram was rerun for two additional seeds on the same fourteen-budget grid. The original and replication seed pairs have high-budget means $25.97$ and $24.79$, respectively, and both produce exactly $35/63$ collapsed two-seed cells. Raw seeded-run collapse counts are $68/126$ and $74/126$. The close reproduction of the aggregate phase diagram strengthens the claim that the one-hop failure landscape is not an accident of one seed pair (Appendix Figure~\ref{fig:k1_replication}). Four-seed comparisons for larger $K$ are not yet complete in this artifact, so we do not pool the new one-hop seeds into the main cross-$K$ table.

\subsection{Secondary diagnostics and their evidential status}

The pre-existing manuscript bundle reports three useful diagnostic patterns: (i) collapsed checkpoints have extreme TD-error tails and critic magnitudes; (ii) three-hop MPI lowers deterministic target-policy displacement relative to one-hop TD3+BC in $88/90$ audited pairs while lowering the final-policy proxy in only $36/90$; and (iii) targeted shadow critics routed through the final actor are less aligned with a simulator return under the first-hop continuation in eight of eight runs. These observations are consistent with critic miscalibration and with the architectural distinction in Proposition~\ref{prop:target_perturbation}.

However, the repository snapshot does not contain the underlying diagnostic logs, checkpoints, or scripts: the root ignore rules exclude logs and parameter files, and the committed figure generator hard-codes aggregate values. We therefore treat these diagnostics as inherited supporting evidence rather than independently reconstructed primary results. In particular, the shadow-critic comparison uses the first-hop continuation as the reference estimand; it does not establish that first-hop routing is better for environment return or that each route is better calibrated to its own continuation.

\subsection{Threats to validity revealed by code audit}
\label{sec:threats}

\paragraph{$K$-dependent critic initialization.}
The released state constructor splits a seed into $K+1$ keys, initializes the $K$ actors from the first $K$ keys, and initializes the critic from the last key. Consequently, the critic key changes with $K$. The main results are paired by integer seed and evaluation episodes, but not by initial critic parameters. Since critic quality is central to the failure mode, this is a material internal-validity threat. A controlled rerun should allocate fixed keys to actor~1 and the critic, then derive auxiliary actor keys separately.

\paragraph{One optimizer step is not a proximal solve.}
Each hop receives one Adam update. Later actors are persistent and independently initialized rather than copied from their references. The theoretical operator in Eq.~\eqref{eq:spi} is therefore not computed exactly, and the conditional descent premise in Proposition~\ref{prop:conditional_descent} is not audited in the available matrices.

\paragraph{The intervention is not compute matched.}
A $K$-hop run performs $K$ actor optimizer calls per delayed update. The original bundle includes small targeted compute-matched controls, but no complete-grid compute-matched comparison. The sweep identifies the implemented bundle, not the isolated value of re-centering or smaller local coefficients.

\paragraph{Scope and uncertainty.}
The cross-$K$ grid uses two seeds on nine locomotion tasks. Environment--budget cells are correlated through common datasets, implementation choices, and reused seed labels. Collapse counts and cellwise comparisons are descriptive summaries rather than independent-sample hypothesis tests. Broader domains, more seeds, and a pre-specified controlled rerun are needed before making a general offline-RL stability claim.
